\section{计算几何}

\subsection{海伦公式}

\[
  p = \frac{a + b + c}{2},\qquad
  S = \sqrt{p\times(p - a)\times(p-b)\times(p-c)}
\]

\subsection{皮克定理}

想要准确使用皮克公式，必须满足以下两个核心条件：

1、多边形的所有顶点必须都在格点（即平面直角坐标系中的整点）上。

2、多边形必须是简单多边形，即它的边不能自我交叉。

$$ A = i + \frac{b}{2} - 1 $$

$A$ ：表示多边形的面积。

$i$ ：表示多边形内部的格点数（internal）。

$b$ ：表示多边形边界上的格点数（boundary）。

\subsection{高斯面积公式}

高斯面积公式（Gauss's Area Formula），又称鞋带公式（Shoelace Formula），是计算平面直角坐标系中简单多边形有向面积的经典算法。

\textbf{数学定义：}

设 $n$ 个顶点按顺序排列为 $(x_1, y_1), \dots, (x_n, y_n)$，多边形面积 $A$ 的计算公式为：

$$ A = \frac{1}{2} \left| \sum_{i=1}^{n} (x_i y_{i+1} - x_{i+1} y_i) \right| $$

其中满足循环规则：$(x_{n+1}, y_{n+1}) = (x_1, y_1)$。

\textbf{几何原理：}

该公式源于二维向量叉积的几何意义。将多边形分解为以原点为公共顶点的 $n$ 个三角形，各三角形的有向面积为对应边向量叉积的一半。累加后，多边形外部的重叠区域因方向相反而相互抵消，最终收敛于多边形的实际几何面积。

\textbf{适用条件与性质：}

\begin{itemize}
  \item 简单多边形约束：多边形边界连续且不能自相交（如蝴蝶结形），否则会导致代数面积求和失真。
  \item 顶点有序性：顶点必须严格按顺时针或逆时针排列，否则无法围成正确的封闭区域。
  \item 有向面积特性：通常规定逆时针排列结果为正，顺时针为负。工程应用中一般取绝对值以获得物理面积。
\end{itemize}

\textbf{复杂度与应用：}

该算法时间复杂度为 $O(n)$，空间复杂度为 $O(1)$，不依赖多边形的凹凸性。因其高效、轻量化的特点，广泛应用于地理信息系统（GIS）、计算机图形学及建筑BIM系统的不规则区域面积测算。
